\documentclass[11pt,a4paper]{article}

\usepackage[utf8]{inputenc}
\usepackage[T1]{fontenc}
\usepackage[ngerman]{babel}
\usepackage{amsmath,amssymb}
\usepackage{geometry}
\usepackage{hyperref}
\usepackage{microtype}

\geometry{margin=2.6cm}
\hypersetup{
  colorlinks=true,
  linkcolor=blue,
  citecolor=blue,
  urlcolor=blue
}

\title{\textbf{Zur arithmetisch-topologischen Neuinterpretation freier Parameter}\\
\large Standardmodell, \texorpdfstring{$\Lambda\mathrm{CDM}$}{Lambda-CDM} und das Bamberger Modell}
\author{Thomas Hoffbauer}
\date{\today}

\begin{document}
\maketitle

\begin{abstract}
Dieser Artikel formuliert einen Berechnungs-Blueprint für die Neuinterpretation der freien
Parameter des Standardmodells der Teilchenphysik und des kosmologischen
$\Lambda\mathrm{CDM}$-Modells im Rahmen der \#Energiedoku. Die zentrale These lautet, dass
die scheinbar empirischen Konstanten nicht fundamental willkürlich sind, sondern als
Einrastpunkte eines ausfrierenden arithmetischen Zahlengitters verstanden werden können.
Riemann-Nullstellen, Hurwitz-Gitter, Selberg-Sieb, Primzahlvierlinge und lokale
Aharonov--Bohm-Phasen bilden dabei die mathematischen Bausteine einer deduktiven
algebraisch-topologischen Geometrie.
\end{abstract}

\section{Motivation}

Das Standardmodell der Teilchenphysik enthält etwa 26 freie Parameter, darunter
Fermionenmassen, Eichkopplungen, Mischungswinkel und Parameter des Higgs-Potentials.
Das kosmologische $\Lambda\mathrm{CDM}$-Modell ergänzt weitere primäre Parameter, etwa
Materiedichten, kosmologische Konstante und das primordiale Spektrum. In der üblichen
Lesart werden diese Größen empirisch bestimmt und anschließend in die Theorie eingesetzt.

Im Bamberger Modell wird diese Rolle umgedeutet. Naturkonstanten erscheinen nicht als
willkürliche Eingabewerte, sondern als mathematisch ausgezeichnete Stabilitätspunkte eines
arithmetischen Raums. Die Grundannahme lautet:
\[
  \text{Naturparameter}
  =
  \text{Einrastpunkte eines ausfrierenden Zahlengitters}.
\]
Die zugrunde liegende Struktur ist der Hurwitz-Raum $\mathcal{H}$, moduliert durch
Riemannsche Flussquanten und lokale Aharonov--Bohm-Phasen.

\section{Das Standardmodell der Teilchenphysik}

\subsection{Fermionenmassen}

In der klassischen Parametrisierung werden die Massen der zwölf Fermionen durch freie
Yukawa-Kopplungen beschrieben. Im Bamberger Modell wird Masse als arithmetische Trägheit
interpretiert: Sie misst die Knotendichte eines Zustands im Gitter gegenüber der Expansion.

Die Massen entstehen als Eigenwerte eines quaternionischen Dirac--Laplace-Operators auf den
diskreten Normschalen des Hurwitz-Gitters:
\[
  N(q), \qquad q\in\mathcal{H}.
\]
Das Elektron bildet das infinitesimale Basisorbital der innersten stabilen Schale
$N(q)=1$. Höhere Generationen, etwa Myon, Tauon und Quarks, entsprechen komplexeren
Schalenüberlagerungen.

Die Massenverhältnisse werden dabei nicht frei gewählt, sondern aus vierdimensionalen
Volumenverhältnissen hyperbolischer Gitterzellen gelesen. Exemplarisch steht die
Proportionalität
\[
  \frac{m_p}{m_e} \approx 6\pi^5.
\]
Lokale Aharonov--Bohm-Phasenverschiebungen
\[
  \Delta\phi_{\mathrm{local}}
\]
wirken als geometrische Dehnungsfaktoren dieser Zellen.

\subsection{Eichkopplungen}

Die drei Eichkopplungen
\[
  g_1,\quad g_2,\quad g_3
\]
beziehungsweise
\[
  \alpha_{\mathrm{QED}},\quad \alpha_w,\quad \alpha_s
\]
beschreiben im Standardmodell die Stärke der Wechselwirkungen. In der arithmetischen
Interpretation sind sie geometrische Reibungskoeffizienten von Wellenfronten an den
Symmetrieachsen des Gitters.

Jede Kopplung wird aus einer differentiellen Winkelabweichung $\delta_k$ berechnet. Diese
Abweichung misst, unter welchem Winkel de-Broglie-Phasenwellen die Symmetrieachsen der
Einheits-Polytope in $\mathcal{H}$ schneiden. Für den elektromagnetischen Fall wird die
Feinstrukturkonstante mit dem Schnittwinkel des 24-Zellers verknüpft:
\[
  \theta_{24} = \frac{3\pi}{8},
  \qquad
  \alpha^{-1} \approx 2\cdot\frac{1}{\delta}.
\]
Die schwache Kopplung $\alpha_w$ wird analog aus Winkeln des 120-Zellers und der
Symmetriegruppe $H_4$ gelesen. Die starke Kopplung $\alpha_s$ folgt aus hochsymmetrischen
Primgitter-Projektionen des $F_4$-Vierbeins.

Das sogenannte Laufen der Kopplungen wird durch die kumulative Dichte der
Riemann-Nullstellen ersetzt. Die Energieabhängigkeit folgt der logarithmischen kinetischen
Zustandsgleichung
\[
  \Delta\phi(\gamma) = a\ln(\gamma)+b.
\]

\subsection{CKM- und PMNS-Mischungswinkel}

Die CKM- und PMNS-Matrizen beschreiben das Mischen von Quark- und Neutrinogenerationen.
Im Bamberger Modell sind Übergänge zwischen Gitterschalen nur über die Multiplikation mit
den 24 Hurwitz-Einheiten möglich. Deshalb erscheinen die Mischungsmatrizen als
Darstellungsmatrizen endlicher quaternionischer Untergruppen.

Die Mischungswinkel sind dadurch keine freien kontinuierlichen Parameter, sondern diskrete
rationale Bruchteile von $\pi$, deformiert durch die asymptotische Phase des
Chebyshev-Bias. Symbolisch:
\[
  \Theta_{\mathrm{mix}}
  \in
  \pi\cdot\mathbb{Q}
  +
  \Delta\phi_{\mathrm{Cheb}}.
\]

\subsection{Higgs-Potential}

Der Vakuumerwartungswert
\[
  v \approx 246\,\mathrm{GeV}
\]
wird als kritische Dichte interpretiert. Er markiert den Punkt, an dem ein expandierendes
Gas natürlicher Zahlen gezwungen wird, im Hurwitz-Gitter einen algebraischen Abschluss
unter Addition und Multiplikation auszubilden.

Diese kritische Dichte wird durch den minimalen Radius bestimmt, an dem das Selberg-Sieb
eine global zusammenhängende Perkolation erlaubt. Die Higgs-Masse $m_H$ entspricht der
infinitesimalen Krümmung dieses geometrischen Phasenübergangs erster Ordnung.

\section{Das kosmologische \texorpdfstring{$\Lambda\mathrm{CDM}$}{Lambda-CDM}-Modell}

\subsection{Baryonische und dunkle Materiedichte}

Die Parameter
\[
  \Omega_b,\qquad \Omega_c
\]
werden im Standardmodell der Kosmologie empirisch aus Beobachtungsdaten bestimmt. In der
\#Energiedoku entspringen sie der Sieb-Effizienz des Raums.

Wenn der kontinuierliche Zahlenstrom abkühlt, entscheidet das mathematische Gewicht des
Selberg-Siebs, welche Moden einfrieren und welche fluide bleiben. Die dunkle
Materiedichte $\Omega_c$ ergibt sich aus der kombinatorischen Mächtigkeit halbzahliger
Hurwitz-Knoten, die auf basalen Primnormen einrasten und ein starres Hintergrundskelett
bilden. Die baryonische Dichte $\Omega_b$ entsteht aus den Selberg-Überlebenden, die als
rationales Gas weiterfließen und in den Quantenkäfigen der Primzahlvierlinge kondensieren.

Das Verhältnis
\[
  \frac{\Omega_c}{\Omega_b}\approx 5{,}3
\]
wird so als Quotient zweier geometrischer Dichten gelesen: der Packungsdichte der
Halbzahlen und der asymptotischen Verteilung der Vierlings-Knoten auf kritischen Schalen.

\subsection{Kosmologische Konstante}

Die kosmologische Konstante $\Lambda$ beziehungsweise $\Omega_\Lambda$ wird als permanenter
arithmetischer Quelldruck interpretiert. Ein unaufhörlicher Zustrom natürlicher Zahlen
liefert pro Takt eine feste Emissionsrate. Dieser Zustrom wirkt als konstante expansive
Vakuumenergiedichte im Kausalitätshorizont
\[
  e_t=t.
\]
Damit ist $\Lambda$ proportional zur Emissionsrate dividiert durch das infinitesimale
vierdimensionale Gittervolumen.

\subsection{Skalarer Spektralindex und Fluktuationsamplitude}

Der skalare Spektralindex $n_s$ und die Fluktuationsamplitude $A_s$ beschreiben das
primordiale Spektrum kosmischer Dichtewellen. Im arithmetischen Modell entstehen diese
Dichtefluktuationskeime durch magnetische Flusschläuche der Riemann-Nullstellen.

Die quantisierte Abstoßung der Nullstellen folgt in der Niveauabstandsanalyse dem
GUE-Gesetz des Gaussian Unitary Ensemble. Das kosmische Leistungsspektrum wird daher über
die Fourier-Transformierte der Paar-Korrelationsfunktion der Riemann-Nullstellen modelliert.
Die Abweichung vom skaleninvarianten Harrison--Zel'dovich-Spektrum
\[
  n_s=1
\]
wird mit dem Steigungsfaktor $a$ der thermodynamischen Zustandsgleichung verknüpft:
\[
  n_s = 1-6a.
\]
Für
\[
  a=0{,}001410
\]
ergibt sich
\[
  n_s \approx 0{,}965,
\]
in Übereinstimmung mit der Größenordnung astrophysikalischer Planck-Satellitendaten.

\section{Berechnungsprinzip}

Der gesamte Vorschlag ersetzt physikalische Fit-Parameter durch zahlentheoretische
Operationen. Die schematische Pipeline lautet:
\[
  \text{Zahlenstrom}
  \xrightarrow{\text{Selberg-Sieb}}
  \mathcal{H}\text{-Gitter}
  \xrightarrow{\text{Riemann-Nullstellen}}
  \text{quantisierte Von-Klitzing-Plateaus}.
\]
Das Selberg-Sieb selektiert stabile Moden, die Riemann-Nullstellen erzeugen diskrete
Flussquanten, und der Chebyshev-Bias polarisiert das System chiral. Die Geometrie der
24 Hurwitz-Einheiten bricht die Phasen in charakteristische Verhältnisse, die mit
$\alpha$, $m_p/m_e$ und kosmologischen Dichteparametern identifiziert werden.

\section{Fazit}

Die Neuinterpretation der freien Parameter des Standardmodells und des
$\Lambda\mathrm{CDM}$-Modells führt zu einem deduktiven Bild: Konstanten sind keine
externen Eingaben, sondern interne Stabilitätswerte einer arithmetischen Matrix. Die
zentralen Rollen spielen der Hurwitz-Raum $\mathcal{H}$, Riemann-Nullstellen,
Primzahlvierlinge, Aharonov--Bohm-Phasen und quantisierte Plateau-Strukturen. In dieser
Lesart erscheint das Universum als geschlossene algebraisch-topologische Ordnung, deren
physikalische Zahlenwerte aus reiner Zahlentheorie hervorgehen.

\end{document}
